\section{Technical Background and Related Work}
\subsection{Scripting Platform}
Google Apps Script supports JavaScript through its V8 runtime. It is not a
standard Node.js or browser runtime, so libraries must be evaluated for that
environment rather than assumed compatible~\cite{googleAppsScriptV8}.
This project therefore adopts Apps Script instead of the earlier Python
command-line assumption.

Google's HTML service supports asynchronous server calls through
\texttt{google.script.run}. A form containing a file-input field can pass the
selected file to a server function as a Blob~\cite{googleHtmlCommunication}.
This supports the proposed local-file upload without requiring the instructor
to manually upload an image to Drive first. The chosen OCR method may still
need temporary cloud files; upload convenience does not establish a
storage-free processing pipeline.

These platform references are technical documentation, not experimental research
on OCR accuracy. The following research provides context for selecting and
evaluating an OCR approach; it does not establish performance on Zoom screenshots.

\subsection{Traditional OCR Pipelines}
Smith's conference paper describes the historical Tesseract engine, including
line finding, word recognition, and adaptive classification~\cite{smith2007tesseract}.
It illustrates that OCR involves locating and organizing text as well as
recognizing characters. The paper concerns the engine described in 2007 and
should not be treated as a specification or benchmark of current Tesseract releases.

For this capstone, the design implication is to evaluate the complete extraction
pipeline. A recognizer may read interface text correctly but still produce an
incorrect student roster if controls or role labels are included. This is our
application-specific inference, not a result reported by Smith. Trials should
therefore distinguish missed name regions, character errors, and unwanted
interface text, and compare an uncropped panel with a name-region crop.

\subsection{Neural Sequence Recognition}
Shi, Bai, and Yao propose an architecture combining feature extraction, sequence
modeling, and transcription for scene-text recognition~\cite{shi2015crnn}.
Their approach handles variable-length sequences without explicit character
segmentation and supports recognition without a predefined lexicon. The cited
version is the authors' 2015 arXiv preprint; its evaluations concern scene-text
benchmarks rather than Zoom participant lists.

This motivates testing names of different lengths and spellings instead of
assuming a fixed dictionary will cover a roster. It does not justify automatically
replacing an unfamiliar name with a common word. Human review remains necessary.
Training a new recognition model is outside this project's initial scope; this
research informs evaluation of existing tools.

\subsection{Practical OCR Systems}
Du et al. describe PP-OCR, a lightweight system with text detection, direction
classification, and recognition components~\cite{du2020ppocr}. Their work studies
model capability and size through optimization strategies and ablation experiments.
The cited source is a 2020 arXiv preprint and describes that version of PP-OCR,
not every later PaddleOCR release.

The relevant lesson for this project is to consider processing cost alongside
recognition quality. We will record time to produce a reviewable roster and
manual correction effort, rather than select a tool solely by an unrelated
benchmark score. A Zoom panel has a narrower layout than the broader text-image
settings addressed by this research; local experiments remain necessary.

\subsection{Connecting OCR to JavaScript Automation}
Tesseract.js provides a JavaScript interface to a WebAssembly port of Tesseract
for browsers and Node.js~\cite{tesseractJs}. This is primary software documentation,
not a peer-reviewed accuracy study. It offers a candidate scripting route:
test recognition in the browser-based sidebar and pass extracted text to the
Apps Script workflow. Compatibility, model loading, worker behavior, and latency
must be demonstrated in the actual sidebar before adopting this approach.

Browser or Node.js support does not imply compatibility with server-side
\texttt{Code.gs}; Apps Script's runtime differs from both~\cite{googleAppsScriptV8}.
An alternative architecture is to upload the image for server-orchestrated OCR.
M2 will compare feasible integration paths, including permissions, storage,
cleanup, and any service cost. No engine or deployment path is selected by this
literature review.

\subsection{Research Questions for the Technology Trial}
\begin{enumerate}
\item Does cropping to the name region reduce unwanted interface text without
losing student entries?
\item How do readable screenshot scale, name length, and duplicate display names
affect extraction errors and human correction effort?
\item Which feasible scripting integration produces a correctable roster with
acceptable latency and manageable authorization and data handling?
\end{enumerate}
These questions translate the literature into project-specific experiments.
They are proposed comparisons, not claims of published or measured improvements.

